% BHU PYQ -- Professional Exam Template
\documentclass[11pt,a4paper]{article}

\usepackage[a4paper,top=2.5cm,bottom=2.5cm,left=2.8cm,right=2.8cm,headheight=14pt]{geometry}
\usepackage{amsmath,amssymb}
\usepackage[T1]{fontenc}
\usepackage[utf8]{inputenc}
\usepackage{lmodern}
\usepackage{microtype}
\usepackage{fancyhdr}
\usepackage{enumitem}
\usepackage{listings}
\usepackage{hyperref}
\usepackage{lastpage}
\usepackage{parskip}

\hypersetup{colorlinks=true,urlcolor=black,linkcolor=black,
  pdftitle={CS-403: Programming For Problem Solving -- BHU PYQ}}

\pagestyle{fancy}
\fancyhf{}
\renewcommand{\headrulewidth}{0.4pt}
\renewcommand{\footrulewidth}{0.4pt}
\fancyhead[L]{\small\textit{Banaras Hindu University}}
\fancyhead[R]{\small\textit{CS-403: Programming For Problem Solvin}}
\fancyfoot[C]{\small Page \thepage\ of \pageref{LastPage}}

\lstset{
  basicstyle=\small\ttfamily,
  frame=single,
  breaklines=true,
  columns=flexible,
  keepspaces=true
}

\newlist{parts}{enumerate}{2}
\setlist[parts,1]{label=(\alph*),leftmargin=*,itemsep=4pt,topsep=3pt}
\setlist[parts,2]{label=(\roman*),leftmargin=*,itemsep=2pt,topsep=2pt}
\newcommand{\pts}[1]{\hfill\textit{[#1]}}

\begin{document}
\begin{center}
  {\large\bfseries BANARAS HINDU UNIVERSITY}\\[2pt]
  {\normalsize B.Sc. (Hons.) Semester IV Examination, 2023-24}\\[6pt]
  \rule{0.8\linewidth}{0.5pt}\\[6pt]
  {\large\bfseries Paper: CS-403 --- Programming For Problem Solving}\\[4pt]
  {\normalsize Time: Three hours \hfill Maximum Marks: 70}
\end{center}

\rule{\linewidth}{0.4pt}

\smallskip
\noindent\textit{Instructions: Question 1 is compulsory, attempt any four questions from the remaining siz questions.}

\bigskip

\medskip\noindent\textbf{Question 100.} |B (
CS - 403 : Programming for Problem Solving
| Note: Question 1 is compulsory, attempt any four questions from the remaining siz questions.

\medskip\noindent\textbf{Question 1.} (a) Explain 5 major types of programming languages with examples.\pts{5}
\begin{parts}
  \item What is the difference between the top-down and bottom-up approaches of programming?\pts{5}
  \item What are the steps involved in solving and problem? Explain your answer with an example.\pts{4}
\end{parts}

\medskip\noindent\textbf{Question 2.} (a) Write an algorithm to find the greatest between 3 numbers.\pts{4}
\begin{parts}
  \item Design an algorithm and flowchart to convert a decimal number to a binary number.\pts{6}
  \item Compare and contrast compilers from interpreters.\pts{4}
\end{parts}

\medskip\noindent\textbf{Question 3.} (a) Describe the order of precedence with regards to operators in C language.\pts{5}
? (b) Explain different data types used in C language?\pts{5}
\begin{parts}
  \item Write the steps involved in converting a C source code to executable code.\pts{4}
\end{parts}

\medskip\noindent\textbf{Question 4.} (a) Write C Program to find the factorial of a number using recursion.\pts{6}
\begin{parts}
  \item What will be the outcome of the following conditional statement if the value of variable s is 10? s >=10 \&\&
  S < 25 \&\& S! =12.\pts{4}
  \item Write a call-by-value function to swap to numbers.\pts{4}
\end{parts}

\medskip\noindent\textbf{Question 5.} (a) Explain with example pre-increment(+-+i) and post-increment i++).\pts{4}
\begin{parts}
  \item Difference between formal argument and actual argument?\pts{3}
  \item What are variables and in what way is it different from constants?\pts{4}
  \item Explain call-by-value and call=by-reference.\pts{3}
\end{parts}

\medskip\noindent\textbf{Question 6.} (a) Write the difference between insertion sort and selection sort. 4]
\begin{parts}
  \item Write the pseudo of the quick sort algorithm.\pts{5}
  \item What will be the output of the following C code?\pts{5}
  2 | \#include<stdio.h>
\begin{lstlisting}[language=C]
2 int main()
3 { .
4 inta=1,b=2,c=3,d=4,e;
5 e=ctd=b*a; is
3 printf("%d, %d\n", e, d); ;
7 t
\end{lstlisting}
\end{parts}

\medskip\noindent\textbf{Question 7.} (a) How Insertion Sort works? Write the pseudo of insertion sort. Explain step-by-step to sort an array A[3, 8,
5, 4, 1, 9, -2] in descending order.\pts{7}
\begin{parts}
  \item What is Searching Algorithm? Explain different types of search operations.
  \item Explain step-by-step to search element ’4’ in array AJ3, 8, 5, 4, 1, 9, -2] using binary search algorithm. 7]
\end{parts}

\vfill
\rule{\linewidth}{0.4pt}
\begin{center}\small
\href{https://bhu-exam.vercel.app/}{Click here to visit our website}\\[4pt]\href{https://me-aryan.vercel.app/}{Click here to visit our contributor}\\[4pt]\href{https://quashan-me.vercel.app/}{Click here to visit our contributor}
\end{center}
\end{document}